\DocumentMetadata{
  pdfversion=2.0,
  pdfstandard=ua-2,
  lang=en-US,
  tagging=on,
  tagging-setup={math/setup=mathml-SE,tabsorder=structure}
}
\documentclass[11pt,letterpaper]{article}
\usepackage[margin=0.68in,includefoot]{geometry}
\usepackage{newcomputermodern}
\usepackage{amsmath,amscd,mathtools}
\usepackage{tikz,adjustbox}
\providecommand{\square}{\mdlgwhtsquare}
\usepackage[hidelinks]{hyperref}
\hypersetup{
  pdftitle={Analysis Qualifying Exam, January 2026, Part 1},
  pdfauthor={University of Florida Department of Mathematics},
  pdfsubject={Graduate mathematics examination},
  pdfdisplaydoctitle=true,
  bookmarksopen=true
}
\setcounter{secnumdepth}{0}
\setlength{\parindent}{0pt}
\setlength{\parskip}{0.28em}
\setlength{\emergencystretch}{3em}
\setlength{\leftmargini}{2.15em}
\setlength{\leftmarginii}{2.65em}
\setlength{\leftmarginiii}{3em}
\pagestyle{empty}
\raggedbottom
\allowdisplaybreaks
\NewTaggingSocketPlug{block/list/label}{examproblem}{%
  \tagstructbegin{tag=Lbl}\tagstructend%
  \tagstructbegin{tag=\LBody}%
  \tagstructbegin{tag=\ProblemHeadingTag,title-o={\ProblemHeadingText}}%
  \tagmcbegin{tag=\ProblemHeadingTag,actualtext={\ProblemHeadingText}}%
  #2%
  \tagmcend%
  \tagstructend%
}
\newenvironment{examproblems}[2]{%
  \def\ProblemHeadingTag{#2}%
  \AssignTaggingSocketPlug{block/list/label}{examproblem}%
  \begin{enumerate}%
  \setcounter{enumi}{#1}%
  \renewcommand{\labelenumi}{\textbf{\theenumi.}}%
  \setlength{\topsep}{0.35em}%
  \setlength{\partopsep}{0pt}%
  \setlength{\itemsep}{0.48em}%
  \setlength{\parsep}{0.18em}%
}{\end{enumerate}}
\newenvironment{examsubparts}{%
  \AssignTaggingSocketPlug{block/list/label}{default}%
  \begin{enumerate}%
  \setlength{\topsep}{0.18em}%
  \setlength{\partopsep}{0pt}%
  \setlength{\itemsep}{0.16em}%
  \setlength{\parsep}{0.1em}%
}{\end{enumerate}}
\newenvironment{examinstructions}{%
  \par\begingroup\itshape
}{%
  \par\endgroup\vspace{0.18em}
}
\makeatletter
\renewcommand\subsection{\@startsection{subsection}{2}{0pt}%
  {-1.25ex plus -.25ex minus -.1ex}%
  {0.55ex plus .1ex}%
  {\normalfont\large\bfseries}}
\renewcommand{\@maketitle}{%
  \newpage\null\vskip -1em%
  {\footnotesize\bfseries
    \href{https://www.ufl.edu/}{University of Florida}, \href{https://math.ufl.edu/}{Department of Mathematics}\par}%
  \vskip -0.5em%
  \tagstructbegin{tag=H1,title={Analysis Qualifying Exam, January 2026, Part 1}}%
  \tagpdfparaOff%
  \tagmcbegin{tag=H1}%
    {\LARGE\bfseries\href{https://gma.math.ufl.edu/exams/analysis-qualifying/}{Analysis Qualifying Exam}, January 2026, Part 1\par}%
  \tagmcend%
  \tagpdfparaOn%
  \tagstructend%
  \vskip 0.25em%
  \tagmcbegin{artifact}%
    \hrule height 1pt%
  \tagmcend%
  \vskip 1em%
}
\makeatother
\title{Analysis Qualifying Exam, January 2026, Part 1}
\author{}
\date{}
\begin{document}
\maketitle
\thispagestyle{empty}
\begin{examinstructions}
Write solutions in a neat and logical fashion, giving complete reasons for all steps. Do three of four.
\end{examinstructions}
\begin{examproblems}{0}{H2}
\def\ProblemHeadingText{Problem 1.}
\item
Do two of three.
\begin{examsubparts}
\item[\textnormal{(i)}]
Give an example of a Lebesgue measurable set \(E \subseteq [0,1]\) that is closed, has positive measure, but contains no non-trivial interval or explain why no such example exists.
\item[\textnormal{(ii)}]
Is the function \(g:\mathbb{R}\to\mathbb{R}\) defined by \(g(x)=\frac{1}{1+x^2}\) the Hardy–Littlewood maximal function of an \(\mathcal{L}^1(\mathbb{R})\) function~\(f\)?
\item[\textnormal{(iii)}]
Does there exist a measure space \((X,\mathcal{M},\mu)\) and a nested decreasing sequence \((E_n)_{n=1}^{\infty}\) of measurable sets such that
\[
\lim_{n\to\infty}\mu(E_n)\ne\mu\left(\bigcap_{n=1}^{\infty}E_n\right)?
\]
\end{examsubparts}
\def\ProblemHeadingText{Problem 2.}
\item
Suppose \(f:[0,1]\times[0,1]\to\mathbb{R}\). Show, if (a)–(c), then \(F:[0,1]\to\mathbb{R}\) defined by
\[
F(x)=\int_{[0,1]}f(x,y)\,d\lambda(y)
\]
 is continuous. Here~\(\lambda\) denotes Lebesgue measure on~\(\mathbb{R}\).
\begin{examsubparts}
\item[\textnormal{(a)}]
for each \(x\in[0,1]\) the function \(g_x(y)=f(x,y)\) is continuous;
\item[\textnormal{(b)}]
for each \(y\in[0,1]\) the function \(h_y(x)=f(x,y)\) is continuous; and
\item[\textnormal{(c)}]
\(f\) is bounded,
\end{examsubparts}
\def\ProblemHeadingText{Problem 3.}
\item
Let \((X,\mathcal{M})\) denote a measurable space and let \(P(X)\) denote the power set of~\(X\). Prove, if \(\mathcal{M}\ne P(X)\) and \(\{x\}\in\mathcal{M}\) for each \(x\in X\), then there exists an \(A\notin\mathcal{M}\otimes\mathcal{M}\) such that \([A]_x,[A]^y\in\mathcal{M}\) for all \(x\in X\). Suggestion: First show \(\delta:X\to X\times X\) defined by \(\delta(x)=(x,x)\) is measurable from \((X,\mathcal{M})\) to \((X\times X,\mathcal{M}\otimes\mathcal{M})\).
\def\ProblemHeadingText{Problem 4.}
\item
Suppose \((X,\mathcal{M},\mu)\) is a measure space and \(f,g:X\to[0,\infty)\) are measurable. Show
\begin{examsubparts}
\item[\textnormal{(a)}]
\(\mu_f:\mathcal{M}\to[0,\infty]\) defined by \(\mu_f(A)=\int \chi_Af\,d\mu=\int_A f\,d\mu\) is a measure; and
\item[\textnormal{(b)}]
\(\int g\,d\mu_f=\int gf\,d\mu\).
\end{examsubparts}
\end{examproblems}
\end{document}
